\chapter{System Design and Architecture}

\section{Introduction}

This chapter presents the design and architectural foundations of the Sahhty platform. Building on the requirements and specifications defined in the previous chapter, it translates functional and non-functional requirements into a concrete technical architecture and a set of structural and behavioral models.

The chapter is organized as follows. We first introduce the global architectural model of the system, describing the three-tier organization that separates the mobile frontend, the API backend, and the persistence layer. We then present the static design of the system through a class diagram that captures the core domain entities and their relationships. Finally, we present the dynamic behavior of the system through sequence diagrams that describe the most critical interaction flows.

\section{Modeling Approach}

\subsection{Overview}

To design a software system, engineers rely on modeling languages that provide a standardized visual vocabulary for describing both the structure and the behavior of a system. In this project, we adopt \textbf{UML (Unified Modeling Language)}, which is the most widely used standard for software system modeling.

UML offers a rich set of diagram types covering multiple perspectives:
\begin{itemize}
    \item \textbf{Structural diagrams} describe what the system is made of: its classes, objects, and their static relationships.
    \item \textbf{Behavioral diagrams} describe how the system works: the sequences of interactions, the states objects can be in, and the activities that take place.
\end{itemize}

In this chapter, we focus on two complementary diagrams:
\begin{enumerate}
    \item The \textbf{Class Diagram}, which models the static structure of the domain.
    \item The \textbf{Sequence Diagram}, which models the dynamic interactions for selected use cases.
\end{enumerate}

This combination allows us to reason about both the data model and the behavior of the system before writing a single line of code, reducing the risk of structural inconsistencies and facilitating communication within the development team.

\subsection{Why UML for Sahhty}

The Sahhty platform is a mobile health application with a clearly defined domain model centered around patients, doctors, appointments, health measurements, and pregnancy tracking. This makes UML particularly well-suited:

\begin{itemize}
    \item The domain has well-defined entities with clear relationships, making the class diagram a natural fit for capturing the data model.
    \item The system involves multi-step user flows (registration, authentication, measurement submission, appointment booking) that benefit from sequence diagrams to verify the correctness of the interaction logic before implementation.
    \item UML diagrams serve as a communication bridge between the frontend and backend teams, ensuring that both sides share the same understanding of data contracts and interaction protocols.
\end{itemize}

\section{System Architecture}

\subsection{Overview}

The Sahhty platform follows a \textbf{three-tier client-server architecture}, a classical and proven pattern for mobile applications backed by a REST API. This architecture separates the system into three independent layers, each with a clearly defined responsibility.

\begin{table}[H]
\centering
\caption{Three-tier architecture of the Sahhty platform}
\renewcommand{\arraystretch}{1.4}
\begin{tabular}{|p{3cm}|p{4cm}|p{6cm}|}
\hline
\textbf{Tier} & \textbf{Technology} & \textbf{Responsibility} \\
\hline
Presentation & Flutter (Dart) & Mobile UI, state management, user interaction, HTTP client \\
\hline
Application & Django + DRF & Business logic, REST API, authentication, AI engine, push notifications, WebSocket \\
\hline
Data & PostgreSQL & Persistent storage of all domain data \\
\hline
\end{tabular}
\end{table}

\begin{figure}[H]
    \centering
    \includegraphics[width=0.85\textwidth]{images/architecture_diagram.png}
    \caption{Global three-tier architecture of the Sahhty platform}
    \label{fig:arch_global}
\end{figure}

\subsection{Presentation Tier}

The presentation tier is the Flutter mobile application targeting Android. It is responsible for all user interface rendering, user input handling, and local state management. The application uses the \textbf{Riverpod} state management library to manage global application state and the \textbf{Dio} HTTP client to communicate with the backend API.

The Flutter application is structured into feature modules, each containing its own screens, providers, and local models. The main modules are:

\begin{itemize}
    \item \texttt{auth} — registration, login, OTP verification, and two-factor authentication.
    \item \texttt{home} — patient dashboard and pregnancy tracking summary.
    \item \texttt{measurements} — health measurement list, add measurement, and AI risk alert dialogs.
    \item \texttt{appointments} — appointment booking, doctor search, and appointment management.
    \item \texttt{alerts} — alert feed with filtering and pagination.
    \item \texttt{medications} — treatment list, DCI search, medication safety analysis.
    \item \texttt{medical\_files} — file upload, in-app viewer, and access management.
    \item \texttt{doctors} (doctor-side) — dashboard, appointment confirmation, schedule management.
    \item \texttt{settings} — language selection, 2FA toggle, profile editing.
\end{itemize}

Navigation between modules is handled declaratively using \textbf{Go Router}, which supports role-based redirects (patient vs doctor home) enforced by a route guard that reads the authentication state from Riverpod.

\subsection{Application Tier}

The application tier is the Django backend, organized as a collection of \textbf{domain applications} each responsible for a specific business area. The backend exposes a \textbf{RESTful API} built with Django REST Framework (DRF) and authenticates all protected requests using \textbf{JWT tokens} issued by Simple JWT.

The backend applications and their API prefixes are:

\begin{table}[H]
\centering
\caption{Django backend applications and their API prefixes}
\renewcommand{\arraystretch}{1.3}
\begin{tabular}{|p{4.5cm}|p{4cm}|p{5cm}|}
\hline
\textbf{Application} & \textbf{URL Prefix} & \textbf{Responsibility} \\
\hline
\texttt{users} & \texttt{/users/} & Authentication, OTP, 2FA, FCM device registration \\
\hline
\texttt{patients} & \texttt{/patients/} & Patient profile management \\
\hline
\texttt{doctors} & \texttt{/doctors/} & Doctor profile, schedule, location \\
\hline
\texttt{appointments} & \texttt{/appointments/} & Booking, confirmation, cancellation \\
\hline
\texttt{measurements} & \texttt{/measurements/} & Health measurements and AI risk assessment \\
\hline
\texttt{alerts} & \texttt{/alerts/} & Alert generation and retrieval \\
\hline
\texttt{Pregnancies} & \texttt{/pregnancies/} & Pregnancy tracking \\
\hline
\texttt{medications} & \texttt{/medications/} & Medication search and comparison \\
\hline
\texttt{dci} & \texttt{/dci/} & DCI database search \\
\hline
\texttt{medical\_files} & \texttt{/medical\_files/} & File upload and access control \\
\hline
\end{tabular}
\end{table}

In addition to the REST API, the backend integrates several cross-cutting services:
\begin{itemize}
    \item \textbf{ML Engine (XGBoost)} — invoked synchronously after each new measurement submission to compute the pregnancy risk level.
    \item \textbf{Firebase Cloud Messaging (FCM)} — used to deliver push notifications to mobile devices for appointment events and high-risk health alerts.
    \item \textbf{Django Channels (WebSocket)} — provides real-time notifications for appointment bookings and medical file access requests.
    \item \textbf{APScheduler} — runs background jobs for periodic alerts (medication reminders, missing measurement alerts, appointment reminders).
    \item \textbf{Brevo (SMTP)} — delivers OTP verification codes and account notifications by email.
\end{itemize}

\subsection{Data Tier}

The data tier is a \textbf{PostgreSQL} relational database. All application data is persisted through Django's ORM, which maps Python model classes to relational tables. PostgreSQL was chosen for its strong support of relational constraints, ACID transactions, and complex queries required by the domain (e.g., joining patient profiles with their measurements and risk assessments in a single query).

\subsection{Communication Protocols}

\begin{table}[H]
\centering
\caption{Communication protocols used in the Sahhty platform}
\renewcommand{\arraystretch}{1.3}
\begin{tabular}{|p{4cm}|p{4cm}|p{5.5cm}|}
\hline
\textbf{From} & \textbf{To} & \textbf{Protocol / Technology} \\
\hline
Flutter App & Django Backend & HTTPS REST (JSON), JWT Bearer Token \\
\hline
Flutter App & Django Channels & WebSocket (WSS), JWT query parameter \\
\hline
Wear OS Watch & Flutter App & Wearable Data Layer API (Bluetooth) \\
\hline
Django Backend & Mobile Device & Firebase Cloud Messaging (FCM) \\
\hline
Django Backend & User Email & SMTP via Brevo \\
\hline
\end{tabular}
\end{table}

\section{Static Design: Class Diagram}

\subsection{Overview}

The class diagram represents the structural backbone of the Sahhty domain. It describes the main entities, their attributes, and the associations between them. The diagram follows standard UML notation and maps directly to the Django ORM model classes implemented in the backend.

The diagram covers the following core domain classes:

\begin{itemize}
    \item \textbf{User} — the base account entity shared by both patients and doctors.
    \item \textbf{Patient} — extends User with medical profile data.
    \item \textbf{Doctor} — extends User with professional profile data.
    \item \textbf{MenstrualCycle} — tracks the menstrual status of a female patient.
    \item \textbf{Pregnancy} — represents an active or past pregnancy linked to a patient.
    \item \textbf{Measurement} — a single health reading (blood pressure, heart rate, etc.).
    \item \textbf{RiskAssessment} — the AI-computed risk evaluation for a patient.
    \item \textbf{Appointment} — a scheduled consultation between a patient and a doctor.
    \item \textbf{DoctorSchedule} — the weekly availability slots defined by a doctor.
    \item \textbf{Alert} — a health or system notification generated by the backend.
    \item \textbf{Treatment} — an active medication plan prescribed for a patient.
    \item \textbf{MedicalFile} — a personal health document uploaded by a patient.
    \item \textbf{MedicalFileAccess} — an access grant allowing a doctor to view a patient's files.
    \item \textbf{FCMDevice} — the push notification token registered by the user's device.
\end{itemize}

\begin{figure}[H]
    \centering
    \includegraphics[width=1.0\textwidth]{images/class_diagram.png}
    \caption{Class diagram of the Sahhty domain model}
    \label{fig:class_diagram}
\end{figure}

\subsection{Key Relationships}

The following table summarizes the most important associations in the domain model:

\begin{table}[H]
\centering
\caption{Key associations in the Sahhty class diagram}
\renewcommand{\arraystretch}{1.4}
\begin{tabular}{|p{3.5cm}|p{3.5cm}|p{2.5cm}|p{4cm}|}
\hline
\textbf{Source} & \textbf{Target} & \textbf{Multiplicity} & \textbf{Semantic} \\
\hline
User & Patient & 1 -- 0..1 & A user may have one patient profile \\
\hline
User & Doctor & 1 -- 0..1 & A user may have one doctor profile \\
\hline
Patient & MenstrualCycle & 1 -- 0..1 & A patient has at most one menstrual cycle record \\
\hline
Patient & Pregnancy & 1 -- 0..* & A patient may have multiple pregnancies \\
\hline
Patient & Measurement & 1 -- 0..* & A patient may have many health measurements \\
\hline
Patient & RiskAssessment & 1 -- 0..* & Multiple risk evaluations over time \\
\hline
Patient & Alert & 1 -- 0..* & Alerts are addressed to a specific patient \\
\hline
Patient & Treatment & 1 -- 0..* & A patient may have multiple active treatments \\
\hline
Patient & MedicalFile & 1 -- 0..* & A patient owns their own medical files \\
\hline
Doctor & DoctorSchedule & 1 -- 0..7 & Up to one schedule entry per day of the week \\
\hline
Appointment & Patient & 0..* -- 1 & Each appointment is associated with one patient \\
\hline
Appointment & Doctor & 0..* -- 1 & Each appointment is associated with one doctor \\
\hline
MedicalFileAccess & Patient & 0..* -- 1 & Access grants reference the patient who owns the files \\
\hline
MedicalFileAccess & Doctor & 0..* -- 1 & Access grants reference the doctor who was granted access \\
\hline
User & FCMDevice & 1 -- 0..* & A user may register multiple devices \\
\hline
\end{tabular}
\end{table}

\subsection{User Entity Design}

The \texttt{User} entity is the central authentication entity of the platform. It extends Django's \texttt{AbstractBaseUser} to replace the default username-based authentication with email-based authentication. The \texttt{role} field (Patient / Doctor) is stored directly on the User entity and determines which profile table is populated after registration.

The User entity incorporates several security-oriented fields:
\begin{itemize}
    \item \texttt{verification\_code} and \texttt{expiration\_date} — for email OTP verification during registration.
    \item \texttt{two\_factor\_code} and \texttt{two\_factor\_enabled} — for optional TOTP-based two-factor authentication.
    \item \texttt{failed\_login\_attempts} and \texttt{lockout\_until} — for brute-force login protection.
    \item \texttt{can\_reset\_password} — a server-side gate for the password reset flow.
\end{itemize}

\subsection{Measurement and Risk Assessment Design}

The \texttt{Measurement} entity uses a flexible dual-value design (\texttt{value1}, \texttt{value2}) to accommodate measurement types that require a single number (heart rate, temperature) and types that require a pair of numbers (blood pressure: systolic and diastolic). The \texttt{context} field distinguishes between measurements entered manually, measurements submitted automatically by the Wear OS watch (\textit{smartwatch data}), and any future data source.

The \texttt{RiskAssessment} entity stores the full feature vector used by the XGBoost model at the time of prediction, providing full auditability of every risk evaluation. It records the global risk level computed by the model, an optional personal risk note based on the patient's known allergies and chronic conditions, and the final combined risk level presented to the user.

\section{Dynamic Design: Sequence Diagrams}

\subsection{Overview}

Sequence diagrams describe the temporal ordering of messages exchanged between the actors and the components of the system to accomplish a specific use case. We present sequence diagrams for the following key scenarios:

\begin{enumerate}
    \item Patient registration with OTP email verification.
    \item Patient login with optional two-factor authentication.
    \item Health measurement submission with AI risk assessment.
    \item Appointment booking by a patient.
    \item Real-time medical file access request via WebSocket.
\end{enumerate}

\subsection{Sequence Diagram 1: Patient Registration}

The registration flow involves the Flutter frontend, the Django Users application, and the Brevo email service. The sequence unfolds as follows:

\begin{enumerate}
    \item The patient fills in the registration form (personal and medical information) and submits it.
    \item The Flutter app sends a \texttt{POST /users/auth/signup/} request to the backend.
    \item The backend validates the submitted data, creates the \texttt{User} and \texttt{Patient} records in the database with \texttt{is\_verified = false}, generates a random 6-digit OTP, and requests Brevo to send it to the patient's email.
    \item The backend returns \texttt{201 Created}.
    \item The Flutter app navigates to the OTP verification screen.
    \item The patient enters the code and submits it via \texttt{POST /users/auth/verify-email/}.
    \item The backend validates the code and its expiration date, sets \texttt{is\_verified = true}, and returns \texttt{200 OK}.
    \item The Flutter app navigates to the login screen.
\end{enumerate}

\begin{figure}[H]
    \centering
    \includegraphics[width=0.90\textwidth]{images/seq_registration.png}
    \caption{Sequence diagram --- Patient registration with OTP verification}
    \label{fig:seq_reg}
\end{figure}

\subsection{Sequence Diagram 2: Patient Login with Two-Factor Authentication}

\begin{enumerate}
    \item The patient submits their email and password via \texttt{POST /users/auth/signin/}.
    \item The backend validates the credentials and checks \texttt{two\_factor\_enabled}.
    \item If 2FA is disabled, the backend returns a JWT access and refresh token pair immediately.
    \item If 2FA is enabled, the backend generates a 6-digit TOTP code, sends it by email via Brevo, and returns a \texttt{200} response with \texttt{"requires\_2fa": true} and no token.
    \item The Flutter app navigates to the 2FA verification screen.
    \item The patient enters the code and submits it via \texttt{POST /users/auth/verify-2fa/}.
    \item The backend validates the code and returns the JWT token pair.
    \item The Flutter app stores the tokens in Flutter Secure Storage, parses the JWT payload to extract the user role, and navigates to the appropriate home screen (patient or doctor).
\end{enumerate}

\begin{figure}[H]
    \centering
    \includegraphics[width=0.90\textwidth]{images/seq_login.png}
    \caption{Sequence diagram --- Patient login with two-factor authentication}
    \label{fig:seq_login}
\end{figure}

\subsection{Sequence Diagram 3: Health Measurement Submission with AI Risk Assessment}

This is the most technically significant flow in the system, as it integrates the REST API layer, the PostgreSQL database, and the XGBoost ML engine in a single synchronous pipeline.

\begin{enumerate}
    \item The patient fills in the measurement form and submits it.
    \item The Flutter app sends a \texttt{POST /measurements/MeasurementService/create\_measurement/} request with the measurement payload (type, value1, value2, unit, context).
    \item The backend saves the \texttt{Measurement} record to the database.
    \item The backend calls \texttt{predict\_risk(patient\_id)} in the ML Engine module.
    \item The ML Engine queries the database for the patient's latest measurements across all types and assembles the feature vector.
    \item The XGBoost classifier processes the feature vector and returns a risk class and a probability score.
    \item The ML Engine converts the probability to a risk percentage and returns \texttt{(risk\_level, risk\_percentage)} to the measurement service.
    \item A \texttt{RiskAssessment} record is created in the database with the computed values.
    \item If the risk level is \texttt{HIGH}, the backend creates an \texttt{Alert} record and sends a push notification via FCM.
    \item The API returns \texttt{201 Created} with the created measurement and the risk fields.
    \item The Flutter app reads \texttt{risk\_level}: if \texttt{medium}, it shows a warning dialog in orange; if \texttt{high}, it shows a critical red dialog with a recommendation to consult a doctor.
\end{enumerate}

\begin{figure}[H]
    \centering
    \includegraphics[width=0.90\textwidth]{images/seq_measurement.png}
    \caption{Sequence diagram --- Measurement submission with AI risk assessment}
    \label{fig:seq_measurement}
\end{figure}

\subsection{Sequence Diagram 4: Appointment Booking by a Patient}

\begin{enumerate}
    \item The patient browses the doctor list and opens a doctor's profile.
    \item The Flutter app fetches the doctor's schedule via \texttt{GET /doctors/DoctorService/\{id\}/schedule/}.
    \item The patient selects an available date and confirms the booking.
    \item The Flutter app sends a \texttt{POST /appointments/AppointmentService/create\_appointment/} request.
    \item The backend creates the \texttt{Appointment} record with status \texttt{PENDING} and triggers a WebSocket notification to the doctor's connected channel.
    \item The backend returns \texttt{201 Created} to the Flutter app.
    \item The Flutter app shows a confirmation dialog and refreshes the appointment list.
    \item On the doctor's device, the WebSocket message is received by the Flutter app in real time and a local notification is displayed.
\end{enumerate}

\begin{figure}[H]
    \centering
    \includegraphics[width=0.90\textwidth]{images/seq_appointment.png}
    \caption{Sequence diagram --- Appointment booking with real-time WebSocket notification}
    \label{fig:seq_appt}
\end{figure}

\subsection{Sequence Diagram 5: Medical File Access Request}

\begin{enumerate}
    \item A doctor searches for a patient by name in the medical file access screen.
    \item The Flutter app calls \texttt{GET /medical\_files/MedicalFileService/search\_patients/} with the search query.
    \item The backend returns the list of matching patients.
    \item The doctor selects a patient and sends an access request via \texttt{POST /medical\_files/MedicalFileService/request\_access/}.
    \item The backend creates a \texttt{MedicalFileAccess} record with status \texttt{PENDING} and sends a WebSocket notification to the patient's connected channel.
    \item On the patient's device, the Flutter WebSocket listener receives the message and displays a real-time notification regardless of which screen the patient is on.
    \item The patient opens the access management screen, which calls \texttt{GET /medical\_files/MedicalFileService/get\_patient\_doctors\_requests/} to list pending requests.
    \item The patient approves the request via \texttt{POST /medical\_files/MedicalFileService/grant\_access/}.
    \item The backend updates the \texttt{MedicalFileAccess} status to \texttt{GRANTED} and sends a WebSocket confirmation to the doctor's channel.
    \item The doctor's Flutter app receives the confirmation in real time and the patient's files become accessible.
\end{enumerate}

\begin{figure}[H]
    \centering
    \includegraphics[width=0.90\textwidth]{images/seq_medical_access.png}
    \caption{Sequence diagram --- Medical file access request via WebSocket}
    \label{fig:seq_access}
\end{figure}

\section{Database Schema Overview}

The relational database is managed by PostgreSQL and populated through Django's ORM migrations. The following table presents the main tables, their primary purpose, and their key foreign key relationships.

\begin{table}[H]
\centering
\caption{Main database tables of the Sahhty platform}
\renewcommand{\arraystretch}{1.3}
\begin{tabular}{|p{4cm}|p{5cm}|p{5cm}|}
\hline
\textbf{Table} & \textbf{Purpose} & \textbf{Key References} \\
\hline
\texttt{users\_user} & User accounts with authentication fields & — \\
\hline
\texttt{patients\_patient} & Patient medical profile & \texttt{users\_user} \\
\hline
\texttt{patients\_menstrualcycle} & Menstrual cycle data & \texttt{patients\_patient} \\
\hline
\texttt{pregnancies\_pregnancy} & Pregnancy tracking & \texttt{patients\_patient} \\
\hline
\texttt{doctors\_doctor} & Doctor professional profile & \texttt{users\_user} \\
\hline
\texttt{doctors\_doctorschedule} & Weekly availability slots & \texttt{doctors\_doctor} \\
\hline
\texttt{measurements\_measurement} & Health readings & \texttt{patients\_patient} \\
\hline
\texttt{measurements\_riskassessment} & AI risk evaluation results & \texttt{patients\_patient} \\
\hline
\texttt{appointments\_appointment} & Doctor-patient consultations & \texttt{patients\_patient}, \texttt{doctors\_doctor} \\
\hline
\texttt{alerts\_alert} & Health and system notifications & \texttt{users\_user} \\
\hline
\texttt{medications\_treatment} & Active treatment plans & \texttt{patients\_patient} \\
\hline
\texttt{medical\_files\_medicalfile} & Uploaded health documents & \texttt{patients\_patient} \\
\hline
\texttt{medical\_files\_medicalfileaccess} & File access grants & \texttt{patients\_patient}, \texttt{doctors\_doctor} \\
\hline
\texttt{users\_fcmdevice} & Push notification tokens & \texttt{users\_user} \\
\hline
\end{tabular}
\end{table}

\section{Conclusion}

This chapter presented the full design of the Sahhty platform at both the architectural and the model level. The three-tier architecture cleanly separates the Flutter mobile frontend from the Django API backend and the PostgreSQL data layer, ensuring maintainability and a clear assignment of responsibilities across the two development teams.

The class diagram established the structural foundation of the domain, identifying fourteen core entities and their relationships. The design decisions documented here --- the dual-value measurement model, the role-based user entity, and the auditable risk assessment record --- directly guided the Django model implementations described in the following chapter.

The sequence diagrams validated the correctness of the most critical interaction flows before implementation, highlighting the integration points between REST, WebSocket, Firebase Cloud Messaging, and the ML engine. Together, these models form the conceptual backbone of the implementation work presented in the next chapter.
